\documentclass[11pt,a4paper,openright]{ltjsbook}

% macOS Hiragino preset (enable in the configured Docker environment)
\usepackage{hiragino-base}
\usepackage{bookmacro-lua}

\title{テンプレート}
\subtitle{これはテンプレートです}
\author{坪井 一馬}
\date{\today}

\begin{document}

\maketitle

\tableofcontents

\chapter{通常章}

\section{著作権フリーのサンプル文章}

\subsection{サンプル1}

　この砂漠で生まれ、この砂漠で育った。ムハディは今日もウリムを連れて山のふもとに行く。
　ウリムはフタコブラクダで、十八歳になる。ムハディより七歳も年寄りの老いぼれラクダだ。
　山のふもとには観光バスの駐車場があって、運がよければ日に15台もやってくる。中国人、ロシア人、インド人、ドイツ人、フランス人、日本人、韓国人、ブラジル人もやってくる。オーストラリア人とイギリス人もやってくるが、ムハディにはほとんど区別がつかない。アメリカ人はめったにやってこないが、来ればすぐにわかる。日本人とおなじくらい、すぐにわかる。
　山といっても、高さがせいぜい150メートルほどの、丘みたいな突起だ。見渡すかぎりの赤茶けた砂漠のなかに、ぴょこんと突き出ている。山の中腹に洞窟がいくつかあって、薄汚れた絵が壁に描かれている。人々はそれを見に来るのだ。ムハディも何度か見たことがあるが、同級生のナジームの絵のほうがよほどうまい。
　運が悪ければバスは一台もやってこない。
　今日はどうだろうか。
　同級生といったが、ムハディはもう学校には行っていない。勤め先の病院で母がロケット弾の直撃を受けて、右足を失った。父はそのすぐあとに部隊に加わった。だから、四人の幼い弟と妹を養うのはムハディの仕事になった。
　金持ちのナジームがまだ学校に行っていることを、ムハディはうらやましくなんかない。ここの仕事は学校なんかよりいろいろと勉強になる。ここが世界の片隅だってことはムハディにもよくわかっているけれど、観光客たちと話をしていると、逆に自分が世界の一番てっぺんにいるような感覚になることもある。
　バスがやってきた。中国人か？　いや、韓国人か日本人だ。年寄りの団体だ。みなころころと太った腹にウエストポーチを大事そうにくくりつけている。
　女のひとりがムハディにカメラを向けながらいう。
「かわいい坊やだわ。うちの孫と同い年くらいよね。学校には行ってないのかしら」
　日本語だ。
　すかさずムハディはいう。
「帰りにラクダ、乗ってよ。写真撮れるよ。いい思い出、撮れるよ」
　一団が歓声をあげる。砂漠の少年の見事な日本語に驚いているのだ。そして少年の透き通った青い目をのぞきこんで、いう。
「かしこい坊やね。わかったわ。きっとあとでラクダに乗ってあげるわね」
　ムハディはときどき、自分の青く透明な目が世界のすべてを見ているような気になることがある。しかし、自分は世界の片隅にいるのだし、自分はおそらく死ぬまでこの地を離れることはないだろうとも思っている。
　ガイドに連れられて山に登っていく観光客たちの姿を、熱くやけた地面に腹ばいになったフタコブラクダのウリムとともに、ムハディはじっと見送っていた。

\subsection{サンプル2}

　丘の上に住んでいる男は、電報がとどくと街に降りてくる。
　足が悪いらしく、杖をついている。かくん、かくんと、曲がったリズムで、乾ききった道を降りてくる。
　杖でかきたてられた土埃が、男の背中にオーラを作る。
　年のころは、そうさな、五十五ってとこか。
　たまに若い女といっしょのこともある。二十歳ぐらいの、髪の長い、ぴちぴちした、はちきれそうな胸と腰つきの女だ。娘じゃないの、とカミさんはいうが、おれはオンナに違いないと踏んでいる。
　五十五なんてまだまだやれるし、痩せてはいるが、骨が太そうな体格だ。こういうやつが一番強い。
　男は電報を握りしめたまま、おれの店の向かいにある銀行に入っていく。しばらくして出てくると、通りを渡ってこの店にやってくる。決まってアニス酒の水割りを一杯注文し、カウンターの上のテレビを食いいるように見つめる。家にゃテレビもないのかね。
　女がいっしょのときは、店には寄らない。そのまま帰ってしまうか、近所で買い物をして帰る。うちには来ない。
　男があの家に住み始めて以来、この半年、ずっとそんな調子だ。
　一度男がひとりのとき、聞いてみた。
「旦那はどんな仕事をなさってるんですかい」
　こっちが質問をしたことを忘れてしまうくらい時間がたってから、男はぼそっと答えた。
「本をな。書いてる」
「どんな本なんですかい」
　それっきり答えはなかった。それ以来、おれは男に話しかけるのをやめた。
　男が電報を忘れたことがあった。おれは悪いと思いながらも、ほかに客もいなかったし、まあ、盗み見をした。
　スイス銀行の口座番号と「14日、20億ドル、要入金」とあった。
　男が忘れ物に気づいて戻ってきた。おれはあわてて電報を戻した。
「読んだのか」
「いえ、滅相もない」
　サザエの口みたいに奥まった目に見据えられて、おれはもう少しで白状しちまうところだった。
「私も世界に関わっている。こんな地の果てにいてもな」
　男が言った。なんのことだか、おれにはわからなかった。
「おまえもそうであるように。私とおまえは、どこか別のところでもつながっている。見ろ」
　男がテレビを杖の先でさした。
　摩天楼に旅客機が突っこんで爆発するところが映っていた。

\section{セクションチェック}

これはサンプルテキストです．\uline{このように長い下線がある場合は，コマンドに注意しなければ，下線が改行されて続かず，非常に読みにくくなってしまうことがあります}．そこで，\kasen{独自に定義したコマンドを導入することによって，LuaLaTeX専用ではあっても下線を改行させて続けることもできます}．


\appendix

\chapter{付録テスト}

\section{セクションチェック}

これはサンプルテキストです．\uline{このように長い下線がある場合は，コマンドに注意しなければ，下線が改行されて続かず，非常に読みにくくなってしまうことがあります}．そこで，\kasen{独自に定義したコマンドを導入することによって，LuaLaTeX専用ではあっても下線を改行させて続けることもできます}．

\subsection{サンプル1}

　この砂漠で生まれ、この砂漠で育った。ムハディは今日もウリムを連れて山のふもとに行く。
　ウリムはフタコブラクダで、十八歳になる。ムハディより七歳も年寄りの老いぼれラクダだ。
　山のふもとには観光バスの駐車場があって、運がよければ日に15台もやってくる。中国人、ロシア人、インド人、ドイツ人、フランス人、日本人、韓国人、ブラジル人もやってくる。オーストラリア人とイギリス人もやってくるが、ムハディにはほとんど区別がつかない。アメリカ人はめったにやってこないが、来ればすぐにわかる。日本人とおなじくらい、すぐにわかる。
　山といっても、高さがせいぜい150メートルほどの、丘みたいな突起だ。見渡すかぎりの赤茶けた砂漠のなかに、ぴょこんと突き出ている。山の中腹に洞窟がいくつかあって、薄汚れた絵が壁に描かれている。人々はそれを見に来るのだ。ムハディも何度か見たことがあるが、同級生のナジームの絵のほうがよほどうまい。
　運が悪ければバスは一台もやってこない。
　今日はどうだろうか。
　同級生といったが、ムハディはもう学校には行っていない。勤め先の病院で母がロケット弾の直撃を受けて、右足を失った。父はそのすぐあとに部隊に加わった。だから、四人の幼い弟と妹を養うのはムハディの仕事になった。
　金持ちのナジームがまだ学校に行っていることを、ムハディはうらやましくなんかない。ここの仕事は学校なんかよりいろいろと勉強になる。ここが世界の片隅だってことはムハディにもよくわかっているけれど、観光客たちと話をしていると、逆に自分が世界の一番てっぺんにいるような感覚になることもある。
　バスがやってきた。中国人か？　いや、韓国人か日本人だ。年寄りの団体だ。みなころころと太った腹にウエストポーチを大事そうにくくりつけている。
　女のひとりがムハディにカメラを向けながらいう。
「かわいい坊やだわ。うちの孫と同い年くらいよね。学校には行ってないのかしら」
　日本語だ。
　すかさずムハディはいう。
「帰りにラクダ、乗ってよ。写真撮れるよ。いい思い出、撮れるよ」
　一団が歓声をあげる。砂漠の少年の見事な日本語に驚いているのだ。そして少年の透き通った青い目をのぞきこんで、いう。
「かしこい坊やね。わかったわ。きっとあとでラクダに乗ってあげるわね」
　ムハディはときどき、自分の青く透明な目が世界のすべてを見ているような気になることがある。しかし、自分は世界の片隅にいるのだし、自分はおそらく死ぬまでこの地を離れることはないだろうとも思っている。
　ガイドに連れられて山に登っていく観光客たちの姿を、熱くやけた地面に腹ばいになったフタコブラクダのウリムとともに、ムハディはじっと見送っていた。

\subsection{サンプル2}

　丘の上に住んでいる男は、電報がとどくと街に降りてくる。
　足が悪いらしく、杖をついている。かくん、かくんと、曲がったリズムで、乾ききった道を降りてくる。
　杖でかきたてられた土埃が、男の背中にオーラを作る。
　年のころは、そうさな、五十五ってとこか。
　たまに若い女といっしょのこともある。二十歳ぐらいの、髪の長い、ぴちぴちした、はちきれそうな胸と腰つきの女だ。娘じゃないの、とカミさんはいうが、おれはオンナに違いないと踏んでいる。
　五十五なんてまだまだやれるし、痩せてはいるが、骨が太そうな体格だ。こういうやつが一番強い。
　男は電報を握りしめたまま、おれの店の向かいにある銀行に入っていく。しばらくして出てくると、通りを渡ってこの店にやってくる。決まってアニス酒の水割りを一杯注文し、カウンターの上のテレビを食いいるように見つめる。家にゃテレビもないのかね。
　女がいっしょのときは、店には寄らない。そのまま帰ってしまうか、近所で買い物をして帰る。うちには来ない。
　男があの家に住み始めて以来、この半年、ずっとそんな調子だ。
　一度男がひとりのとき、聞いてみた。
「旦那はどんな仕事をなさってるんですかい」
　こっちが質問をしたことを忘れてしまうくらい時間がたってから、男はぼそっと答えた。
「本をな。書いてる」
「どんな本なんですかい」
　それっきり答えはなかった。それ以来、おれは男に話しかけるのをやめた。
　男が電報を忘れたことがあった。おれは悪いと思いながらも、ほかに客もいなかったし、まあ、盗み見をした。
　スイス銀行の口座番号と「14日、20億ドル、要入金」とあった。
　男が忘れ物に気づいて戻ってきた。おれはあわてて電報を戻した。
「読んだのか」
「いえ、滅相もない」
　サザエの口みたいに奥まった目に見据えられて、おれはもう少しで白状しちまうところだった。
「私も世界に関わっている。こんな地の果てにいてもな」
　男が言った。なんのことだか、おれにはわからなかった。
「おまえもそうであるように。私とおまえは、どこか別のところでもつながっている。見ろ」
　男がテレビを杖の先でさした。
　摩天楼に旅客機が突っこんで爆発するところが映っていた。

\end{document}